\chapter{Vize}

Vedoucí kuchyně při sestavování jídelníčku potřebuje vědět, kolik porcí
jednotlivých jídel lze připravit ze surovin na skladě. Samotný údaj
o~množství nestačí. Pokud část smetany spotřebuje dnešní oběd a~zbytek
má zítra poslední den spotřeby, ovlivní to výběr jídla na zítřek
i~množství smetany, které bude třeba nakoupit.

Personál kuchyně při příjmu zaznamená množství suroviny a~datum
spotřeby. Stejná surovina může být na skladě několikrát s~různými
termíny spotřeby. Při plánování na konkrétní den vedoucí uvidí, kolik
porcí jednotlivých jídel lze připravit a~které suroviny je vhodné
spotřebovat přednostně.

Do denního plánu vedoucí vybere jídla a~zadá požadované počty porcí.
Pokud se například smetana používá v~polévce i~v~omáčce, nákupní
seznam zohlední spotřebu obou jídel dohromady. Vedoucí tak před nákupem
vidí, kolik smetany skutečně chybí. Po přípravě jídel personál plán
potvrdí a~spotřeba se odepíše nejdříve ze zásob s~nejbližším datem
spotřeby.
